\documentclass[conference]{IEEEtran}
\usepackage{cite}
\usepackage{graphicx}
\usepackage{amsmath}
\usepackage{booktabs}
\usepackage{hyperref}

\begin{document}

\title{Hybrid Retrieval-Augmented Generation for Multi-Paper Research Assistance}

\author{\IEEEauthorblockN{Research RAG Team}
\IEEEauthorblockA{Gen AI / research_rag project}}

\maketitle

\begin{abstract}
We present a hybrid Retrieval-Augmented Generation (RAG) system designed
for multi-paper research assistance. The system combines dense and sparse
text retrieval with a lightweight knowledge graph to improve evidence
retrieval and grounding for downstream LLM-based answer generation.
This paper describes the implementation, experimental protocol used to
compare text-only, graph-only, and hybrid retrieval modes, and provides
the LaTeX source and instructions to reproduce experiments on the
project's artifact dataset.
\end{abstract}

\section{Introduction}
Retrieval-Augmented Generation (RAG) pairs retrieval components with a
language model to ground generation in external documents. In this
work we implement a hybrid index that merges dense (semantic) and
sparse (BM25) text retrieval for document chunks and augments this with
a knowledge graph built from structured extractions. The hybrid design
aims to: (1) improve recall of relevant textual evidence, and (2)
surface structured relations useful for comparative questions.

\section{System Architecture}
Our implementation uses three main components:
\begin{itemize}
  \item \textbf{Chunk Index (text):} Document chunks are embedded with
    a SentenceTransformer model and indexed with FAISS for dense
    retrieval; BM25 (rank\\_bm25) provides sparse lexical scores. A
    weighted combination of dense and sparse scores produces the final
    chunk ranking.
  \item \textbf{Knowledge Graph (graph):} Paper-level extractions are
    converted to a directed graph via \texttt{networkx}. Nodes include
    papers, methods, datasets, metrics, and limitations; edges encode
    relations (e.g., "proposes", "evaluated\_on").
  \item \textbf{Hybrid Retrieval}: The hybrid pipeline performs both
    chunk search and graph node neighborhood retrieval. Retrieved
    items are concatenated into a compact context and passed to an
    LLM prompt template (the project's \texttt{MultiPaperRAGGenerator}).
\end{itemize}

Implementation files of interest: [multi_generator.py](multi_generator.py), [graph_index.py](graph_index.py), [multi_rag.py](multi_rag.py).

\section{HybridIndex Details}
The \texttt{HybridIndex} (in [graph_index.py](graph_index.py)) builds:
\begin{enumerate}
  \item Dense FAISS indices for chunk and node embeddings (cosine via normalized IP).
  \item A BM25 index for lexical matching over chunk tokens.
  \item A combined score: score $= \alpha\cdot s_{dense} + (1-\alpha)\cdot s_{sparse}$.
\end{enumerate}

Graph neighborhoods are retrieved via embedding-based node search and
then expanded with an ego-graph to provide relational context.

\section{Experimental Design}
We evaluate three retrieval modes:
\begin{itemize}
  \item \textbf{Text-only:} Use chunk ranking (dense+sparse) and ignore graph results.
  \item \textbf{Graph-only:} Use node embedding search and neighborhoods only.
  \item \textbf{Hybrid:} Use both text and graph retrieval; consider a ground-truth hit if either source returns evidence.
\end{itemize}

Dataset: we use the project's extracted artifacts (JSON files in \texttt{artifacts/}).

Metrics: For reproducible automated evaluation without invoking an LLM,
we measure retrieval recall@k relative to ground-truth fields in the
artifact (e.g., \texttt{method.core\_idea}, \texttt{evaluation.datasets},
and entries in \texttt{evidence\_map}). For each query derived from a
paper (e.g., "What is the core idea of PAPER?"), we check whether the
ground-truth string appears in the top-k chunk text or graph neighborhood.

\section{Runner and Reproducibility}
We include an experiment runner script at [experiments/compare_rag.py](experiments/compare_rag.py)
that:
\begin{enumerate}
  \item Loads artifact JSONs and synthesizes lightweight text chunks from
    titles, method core ideas, and evidence snippets.
  \item Builds the combined knowledge graph using \texttt{GraphBuilder}.
  \item Constructs a \texttt{HybridIndex} and runs the three retrieval modes.
  \item Reports recall@k aggregated across queries.
\end{enumerate}

To reproduce experiments locally (UNIX):
\begin{verbatim}
python3 -m pip install -r requirements.txt
PYTHONPATH="./research_rag" python3 research_rag/experiments/compare_rag.py
\end{verbatim}

Note: the environment must include the following packages (typical):
\texttt{faiss-cpu}, \texttt{sentence-transformers}, \texttt{rank_bm25}, \texttt{networkx}.
If these are not installed, the runner will fail when building the index.

\section{Results}
The runner reports three aggregate recall measures: text-only, graph-only,
and hybrid recall@k. In this repository a lightweight run requires the
dependencies above; if you run the included script it will print the
summary table. A representative (placeholder) table layout is below —
the numbers are to be filled by running the script in your environment.

\begin{table}[ht]
\caption{Retrieval recall@k (placeholder values)}
\centering
\begin{tabular}{lccc}
\toprule
Mode & Text recall@5 & Graph recall@5 & Hybrid recall@5 \\
\midrule
Text-only & 0.52 & - & 0.52 \\
Graph-only & - & 0.28 & 0.28 \\
Hybrid & 0.52 & 0.28 & 0.64 \\
\bottomrule
\end{tabular}
\end{table}

Replace the above with measured values after running the script.

\section{Discussion}
Preliminary expectations: the hybrid approach should improve overall
recall by capturing evidence that is either lexically/densely similar
in chunks or present as structured nodes in the graph (e.g., dataset
names, limitations). The graph provides structured relations valuable
for comparison-style queries, while dense text retrieves paraphrased
evidence.

Limitations: the evaluation above measures retrieval recall only; a
full evaluation needs LLM-generated answers scored with human
judgements or automatic QA metrics (BLEU/ROUGE/QA-based F1) after
generation. LLM access and rate limits affect end-to-end experiments.

\section{Conclusion}
We described a hybrid RAG implementation for multi-paper research
assistance, provided a reproducible runner to compare retrieval modes,
and included the LaTeX source. Running the provided experiment script
in an environment with the required packages will produce quantitative
recall comparisons; end-to-end LLM evaluation requires valid API keys
and will add grounding/hallucination analysis.

\section*{Acknowledgment}
This work builds on open-source components: FAISS, SentenceTransformers,
networkx and the project's own extractor and graph-builder modules.

\begin{thebibliography}{1}
\bibitem{rag}
S. Lewis et al., "Retrieval-Augmented Generation for Knowledge-Intensive NLP Tasks," 2020.
\bibitem{faiss}
J. Johnson et al., "Billion-scale similarity search with GPUs," 2017.
\end{thebibliography}

\end{document}
